% Generated draft: requires mathematical and source review.
\documentclass[11pt]{article}
\usepackage[margin=25mm]{geometry}
\usepackage{fontspec}
\setmainfont{Latin Modern Roman}
\usepackage{amsmath,amssymb,mathtools,unicode-math}
\setmathfont{Latin Modern Math}
\usepackage{enumitem,needspace,xurl,hyperref}
\setlength{\parindent}{0pt}
\setlength{\parskip}{5pt}
\setlength{\emergencystretch}{3em}
\newcommand{\N}{\mathbb{N}}
\newcommand{\Z}{\mathbb{Z}}
\newcommand{\Q}{\mathbb{Q}}
\newcommand{\R}{\mathbb{R}}
\newcommand{\C}{\mathbb{C}}
\newcommand{\F}{\mathbb{F}}
\newcommand{\A}{\mathbb{A}}
\newcommand{\PP}{\mathbb{P}}
\newcommand{\E}{\mathbb{E}}
\newcommand{\eps}{\varepsilon}
\newcommand{\dd}{\,\mathrm{d}}
\begin{document}
\section*{UnsolvedMath ID: 7200084}
\textbf{Source problem code:} catalogue entry
\subsection{Definitions and mathematical statement}
The source asks how much the endpoint separation of a long radius-one rope decreases when a tight open knot is tied. To distinguish the quantities, for an inextensible rope with centerline \(\gamma:[0,L]\to\mathbb R^3\) parametrized by arclength, define
\[
 d(\gamma)=|\gamma(L)-\gamma(0)|,\qquad
 \Delta(\gamma)=L-d(\gamma).
\]
The reference straight rope has separation \(L\). Radius \(1\) would ordinarily be modeled by a non-self-overlapping unit-thickness tube, with a precise touching/end-cap convention; the source does not supply that model.

A possible tension-tightening formulation, not a source-specified definition, would choose an admissible class \(\mathcal C_K(L)\) of unit-thickness curves of the intended open-knot type \(K\) and ask for
\[
 \Delta_K(L)=L-\sup_{\gamma\in\mathcal C_K(L)}d(\gamma),
\]
or for a sufficiently-long-rope limit of this quantity. The excerpt does not select this optimization, a mechanical equilibrium, or another meaning of ``tight'', and does not state whether a formula may depend on \(K\), \(L\), or additional parameters.

An open arc with freely moving ends does not carry the same fixed knot type as a closed loop. Endpoint constraints, a chosen closure, or a long-knot convention such as straight behavior outside a compact set are needed before \(\mathcal C_K(L)\) is meaningful. The source gives none, and also does not specify stretching, friction, tail directions, or what ``long'' requires. No closed-knot ropelength minimization is silently substituted for this endpoint-shortening question.

The summary reports no verified general expression and points to related tight-knot literature, not a precise theorem answering this question. The pinned Wikipedia open label and search report remain source assessments pending primary-source review.
\subsection{Short English statement}
Find the endpoint shortening caused by a tight open knot in a radius-one rope. A precise answer needs the knot convention, rope length and thickness model, endpoint constraints, and definition of tightening; the source supplies none of these choices.
\subsection{Sources}
\begin{itemize}
\item Local registry: \url{data/unsolvedmath/problems.json}, permanent UnsolvedMath ID \(7200084\), source code catalogue entry; supplied statement, background, and research summary, updated 2026-08-06.
\item Wikipedia contributors, \emph{List of unsolved problems in mathematics}, Geometry bullet 84, pinned revision 1366636767: \url{https://en.wikipedia.org/w/index.php?title=List_of_unsolved_problems_in_mathematics&oldid=1366636767#Geometry}. MediaWiki-source extraction accessed 2026-07-29; primary-source status review remains required.
\item ``Knot Tightening by Constrained Gradient Descent'', arXiv:1002.1723; ``The Shapes of Tight Composite Knots'', arXiv:1110.3262. These are supplied as related literature, without authors or an asserted general shortening formula.
\end{itemize}
\end{document}
